% Generated draft: requires mathematical and source review.
\documentclass[11pt]{article}
\usepackage[margin=25mm]{geometry}
\usepackage{fontspec}
\setmainfont{Latin Modern Roman}
\usepackage{amsmath,amssymb,mathtools,unicode-math}
\setmathfont{Latin Modern Math}
\usepackage{enumitem,needspace,xurl,hyperref}
\setlength{\parindent}{0pt}
\setlength{\parskip}{5pt}
\setlength{\emergencystretch}{3em}
\newcommand{\N}{\mathbb{N}}
\newcommand{\Z}{\mathbb{Z}}
\newcommand{\Q}{\mathbb{Q}}
\newcommand{\R}{\mathbb{R}}
\newcommand{\C}{\mathbb{C}}
\newcommand{\F}{\mathbb{F}}
\newcommand{\A}{\mathbb{A}}
\newcommand{\PP}{\mathbb{P}}
\newcommand{\E}{\mathbb{E}}
\newcommand{\eps}{\varepsilon}
\newcommand{\dd}{\,\mathrm{d}}
\begin{document}
\section*{UnsolvedMath ID: 30002143}
\textbf{Source problem code:} catalogue entry
\subsection{Definitions and mathematical statement}
In the usual smooth interpretation, a contact form on a manifold \(M\) of dimension \(2n+1\) is a one-form \(\alpha\) satisfying \(\alpha\wedge(\dd\alpha)^n\neq 0\) everywhere, and its contact distribution is \(\xi=\ker\alpha\). Define the strict contactomorphism group and the contactomorphism group by
\[
\operatorname{Cont}(M,\alpha)=\{f\in\operatorname{Diff}(M):f^*\alpha=\alpha\},\qquad
\operatorname{Cont}(M,\xi)=\{f\in\operatorname{Diff}(M):\dd f(\xi)=\xi\}.
\]
Here \(\dd f\) denotes the differential of \(f\). The subscript \(0\) denotes the identity component, with the group topology to be specified. Preservation of \(\alpha\) implies preservation of \(\xi\), giving the inclusion
\[
j:\operatorname{Cont}_0(M,\alpha)\hookrightarrow\operatorname{Cont}_0(M,\xi),\qquad j(f)=f.
\]
The supplied question asks what geometric properties this inclusion has. No metric, geometric structure on these groups, or particular property to establish is specified, so this part does not determine a mathematical assertion.

For a symplectic manifold \((X,\omega)\), write \(\operatorname{Ham}(X,\omega)\) for the group of time-one maps of Hamiltonian isotopies, subject to whatever support and boundary conditions are intended. A real-valued quasimorphism on a group \(G\) is a function \(q:G\to\R\) with finite defect
\[
D(q)=\sup_{g,h\in G}|q(gh)-q(g)-q(h)|<\infty.
\]
This supplies a conventional meaning of quasimorphism; the record does not specify homogeneity, continuity, or other regularity requirements.

In the prequantization case, the intended additional data are a principal circle bundle \(\pi:M\to X\) with a connection contact form \(\alpha\) whose curvature corresponds to \(\omega\); the circle and curvature normalizations are not supplied. If these data give the intended projection
\[
p:\operatorname{Cont}_0(M,\alpha)\longrightarrow\operatorname{Ham}(X,\omega),
\]
then an exact extension question for a specified quasimorphism \(q:\operatorname{Ham}(X,\omega)\to\R\) would be whether there exists a quasimorphism \(\widetilde q:\operatorname{Cont}_0(M,\xi)\to\R\) such that
\[
\widetilde q\circ j=q\circ p.
\]
This is a conditional formulation of a possible intended question, not an assertion that the source fixes this interpretation. The record does not specify which quasimorphisms are under consideration, whether extension is exact or only up to bounded error, or whether universal covering groups are intended. It also does not give the dimensions, compactness assumptions, boundary conditions, or Hamiltonian support conventions. Consequently, the supplied record describes a research direction rather than a fully quantified conjecture.
\subsection{Short English statement}
The question concerns the inclusion of strict contactomorphisms into contactomorphisms preserving the contact distribution. In a prequantization setting, it also asks whether Hamiltonian quasimorphisms, pulled back to strict contactomorphisms, extend to the larger group. The intended geometry, hypotheses, and extension quantifiers are unspecified, so no determinate conjecture can be extracted without further source clarification.
\subsection{Sources}
\begin{itemize}
\item Supplied local UnsolvedMath registry record, permanent ID \(30002143\), source problem code \texttt{catalogue entry}, titled ``Geometry and Quasimorphism Extension Along Contactomorphism Inclusions.'' The supplied proposed year is 2012; no individual proposer is identified.
\item Supplied report attribution: \emph{Geometric Group Theory, Hyperbolic Dynamics and Symplectic Geometry} (2013), Oberwolfach Reports, pp.\ 2139--2203, DOI \texttt{10.4171/owr/2012/35}. \url{https://doi.org/10.4171/owr/2012/35}. The supplied background describes Oberwolfach Reports as documenting mathematical workshops and their open-problem discussions.
\item Supplied literature source: \emph{Chern--Weil theory and strict contactomorphisms}. \url{https://arxiv.org/abs/1409.1913}.
\item Supplied literature source: \emph{Oberwolfach Report 35}. \url{https://doi.org/10.4171/owr/2012/35}.
\item Literature assessment dated exactly \texttt{2026-08-21}: the supplied curation labels the record open because it is a context-only follow-up without a determinate geometric claim whose resolution can be verified. This is supplied dated curation derived from the OWR material, not an independently verified assessment of mathematical status. The supplied background calls for independent verification of time-sensitive status and cited claims.
\item The registry labels the statement \texttt{corrected\_verified} and reports comparison with the official Oberwolfach report, with notation, context, and question form normalized without changing the source-level claim. That verification description is supplied provenance and has not been independently checked here.
\end{itemize}
\end{document}
